%% Chapter 5 — Section 5.3: Convergence of the Gibbs Sampler
%% A First Course in Monte Carlo Methods
%% D. Sanz-Alonso & O. Al-Ghattas

\section{Convergence of the Gibbs Sampler}
\label{sec:5.3}

As we know, detailed balance implies that the Gibbs sampler has the desired invariant
distribution. Therefore, if the chain is ergodic, then its distribution converges to the target.
The following lemma gives a sufficient condition for irreducibility.

\begin{lemma}
\label{lem:5.7}
Suppose that the target $f(x_1, \ldots, x_d)$ satisfies:
\[
  \underbrace{f_i(x_i)}_{\text{Marginal}} > 0
  \quad \forall\, i \in \{1, \ldots, d\}
  \implies f(x_1, \ldots, x_d) > 0.
\]
Then, the Gibbs sampler is $f$-irreducible.
\end{lemma}

\begin{example}[Reducible Gibbs Sampler]
\label{ex:5.8}
Consider the following distribution in $\R^2$:
\[
  f(x_1, x_2) = \frac{1}{2}\mathbf{1}_{\{0 \le x_1 \le 1,\, 0 \le x_2 \le 1\}}
    + \frac{1}{2}\mathbf{1}_{\{-1 \le x_1 \le 0,\, -1 \le x_2 \le 0\}},
\]
which is depicted in Figure~\ref{fig:5.2}. The full conditionals are:
\[
  f(x_1|x_2) =
  \begin{cases}
    \mathbf{1}_{\{x_1 \in [0,1]\}} & x_2 \in (0,1),\\
    \mathbf{1}_{\{x_1 \in [-1,0]\}} & x_2 \in (-1,0),
  \end{cases}
\]
\[
  f(x_2|x_1) =
  \begin{cases}
    \mathbf{1}_{\{x_2 \in (0,1)\}} & x_1 \in (0,1),\\
    \mathbf{1}_{\{x_2 \in (-1,0)\}} & x_1 \in (-1,0).
  \end{cases}
\]

\begin{figure}[htbp]
  \centering
  \includegraphics[width=0.55\textwidth]{figures/fig_05_2}
  \caption{\textit{Target density for which the Gibbs sampler is reducible.}}
  \label{fig:5.2}
\end{figure}

If we used a Gibbs sampler, it would be reducible: if we start in one of the two squares, it
would never reach the other square. Note that the condition of the previous lemma fails here:
\begin{align*}
  f_1(x_1) &= \frac{1}{2}\mathbf{1}_{\{-1 \le x_1 \le 1\}},\\
  f_2(x_2) &= \frac{1}{2}\mathbf{1}_{\{-1 \le x_2 \le 1\}},\\
  f(-0.5, 0.5) &= 0 \quad \text{but} \quad f_1(-0.5) > 0, \quad f_2(0.5) > 0.
\end{align*}
\end{example}

\begin{remark}
\label{rem:5.9}
It can be shown that Harris recurrence holds if the transition kernel is absolutely continuous
with respect to an appropriate dominating measure.
\end{remark}

\subsection{Convergence for Finite State Distributions}
\label{sec:5.3.1}

Convergence of the Gibbs sampler in a finite state space can be established using the results on
ergodicity of Markov chains reviewed in Appendix~A. Proposition~5.5 ensures that the target is
the invariant distribution of the chain; convergence to the target follows if we guarantee that
the chain is irreducible and aperiodic.

\begin{theorem}[Convergence of Gibbs Sampler in Finite State Space]
\label{thm:5.10}
Consider a target p.m.f.\ $f(x)$ on a finite state space $E = \{1, \ldots, d\}$, and assume
that $f(x) > 0$ for all $x \in E$. Suppose we run a Gibbs sampler on $f$ by sampling each of
the full conditionals $f_j$ in sequential order. Then, there is $\epsilon > 0$ such that, for
any $x \in E$,
\[
  d_{\mathrm{TV}}\bigl(p_{\mathrm{GS}}^n(x, \cdot),\, f\bigr) \le (1 - \epsilon)^n,
\]
where $p_{\mathrm{GS}}$ denotes the Markov kernel defined by one swipe of sampling all full
conditionals $f_j$, $1 \le j \le d$.
\end{theorem}

\begin{proof}
The assumption that $f(x) > 0$ for all $x \in E$ implies that all full conditionals are
positive. Thus, the Markov kernel $p_{\mathrm{GS}}$ is bounded below by a positive constant and
the result follows from Theorem~A.25 on ergodicity of Markov chains in finite state space.
\end{proof}

\subsection{Convergence for Gaussian Distributions}
\label{sec:5.3.2}

Now we extend our objective to continuous distributions. Instead of investigating a general
continuous distribution, we will focus on multivariate Gaussian targets
$f = \Normal(\mu, \Sigma) = \Normal(\mu, H^{-1})$. Moreover, we restrict our analysis to the
following vanilla Gibbs sampler that makes deterministic block updates:

\begin{algorithm}[H]
\caption{Deterministic Update Gibbs Sampler (DUGS)}
\label{alg:5.2}
\begin{algorithmic}[1]
\Require target $f(x_1, \ldots, x_d)$, initialization
  $X^{(0)} = (X_1^{(0)}, \ldots, X_d^{(0)})$, sample size $N$.
\For{$n = 1, \ldots, N$}
  \For{$j = 1, \ldots, d$}
    \State Sample $X_j^{(n)} \sim f_j\!\bigl(X_j \mid X_1^{(n)}, \ldots, X_{j-1}^{(n)},
      X_{j+1}^{(n-1)}, \ldots, X_d^{(n-1)}\bigr)$.
    \State Update the $j$-th coordinate by setting
      \[
        X^{(n)} = \Bigl(X_1^{(n)}, \ldots, X_{j-1}^{(n)},\, X_j^{(n)},\,
          X_{j+1}^{(n-1)}, \ldots, X_d^{(n-1)}\Bigr).
      \]
  \EndFor
  \State \textbf{end for}
  \State Record sample $X^{(n)}$.
\EndFor
\Ensure Sample $\{X^{(n)}\}_{n=1}^N$.
\end{algorithmic}
\end{algorithm}

We first find a compact way to write the deterministic Gibbs updates. We decompose the precision
matrix $H = H_L + H_D + H_U$ into its lower-triangular, diagonal, and upper-triangular
components. We have the following lemma.

\begin{lemma}
\label{lem:5.11}
Let $\{X^{(n)}\}_{n \ge 0}$ be the Gibbs output with target $f = \Normal(\mu, \Sigma) =
\Normal(\mu, H^{-1})$ and initial state $X^{(0)}$. Then, $\{X^{(n)}\}_{n \ge 0}$ satisfies
the following recurrence:
\begin{equation}
  X^{(n+1)} - \mu = -(H_L + H_D)^{-1} H_U (X^{(n)} - \mu) + (H_L + H_D)^{-1} U^{(n)},
  \label{eq:5.3}
\end{equation}
where $U^{(n)} \sim \Normal(0, H_D)$.
\end{lemma}

\begin{proof}
Denote the state at iteration $n$ by $x = (x_1, x_2, \ldots, x_d)$ and the state at iteration
$n+1$, after a Gibbs update of all coordinates, by $x' = (x_1', x_2', \ldots, x_d')$. Using
Example~\ref{ex:5.2} we can write the Gibbs update on block $j$ as
\[
  H_{jj}(x_j' - \nu_j) = u_j,
\]
where $u_j \sim \Normal(0, H_{jj})$ with
\[
  \nu_j = \mu_j - H_{jj}^{-1}\sum_{i < j} H_{ji}(x_i' - \mu_i)
    - H_{jj}^{-1}\sum_{i > j} H_{ji}(x_i - \mu_i)
\]
since we are conditioning on $x_1', x_2', \ldots, x_{j-1}', x_{j+1}, \ldots, x_d$.
Expanding out $\nu_j$ gives
\[
  H_{jj}(x_j' - \mu_j)
    + [H_{j1};\ldots;H_{j,j-1}]^T(x_{1:j-1}' - \mu_{1:j-1})
    + [H_{j,j+1};\ldots;H_{jd}]^T(x_{j+1:d} - \mu_{j+1:d})
  = u_j.
\]
Therefore, a full Gibbs update can be written in vectorized form as
\[
  (H_L + H_D)(x' - \mu) + H_U(x - \mu) = u,
\]
where $u \sim \Normal(0, H_D)$ as desired.
\end{proof}

The matrix
\[
  B = -(H_L + H_D)^{-1} H_U
\]
is known as the Gauss--Seidel companion matrix. The Gauss--Seidel recurrence in~\eqref{eq:5.3}
belongs to a larger family of iterative techniques to solve linear systems of equations via
matrix splitting. It is clear that a sufficient condition for convergence is that $\rho(B) < 1$.
You will show in Exercise~5.6 that this condition holds provided that the matrix $H$ is
symmetric and positive definite.

Recall from Exercise~3.7 in Chapter~3 that the $\chi^2$-divergence between Gaussian
distributions $f = \Normal(\mu, \Sigma)$ and $\tilde{f} = \Normal(\tilde{\mu}, \tilde{\Sigma})$
is given by
\begin{equation}
  d_{\chi^2}(f \| \tilde{f})
  = \frac{|W|}{\sqrt{|2W - I|}}\,e^{u^T(2W-I)^{-1}\Sigma^{-1}u} - 1
  \label{eq:5.4}
\end{equation}
with $W = \tilde{\Sigma}\Sigma^{-1}$ and $u = \mu - \tilde{\mu}$.

\begin{theorem}[Convergence of Gibbs Sampler for Gaussian Targets]
\label{thm:5.12}
Let the target distribution of the Deterministic Update Gibbs Sampler (DUGS) be
$f = \Normal(\mu, H^{-1}) = \Normal(\mu, \Sigma)$, and let $p_{\mathrm{DUGS}}$ denote the
Markov kernel defined by one full swipe of the DUGS scheme. Then,
\[
  \Expect_f\!\Bigl[d_{\chi^2}\bigl(p_{\mathrm{DUGS}}^n(x^{(0)}, \cdot)\,\|\, f\bigr)\Bigr]
  \sim o\bigl(\rho(B)^n\bigr),
\]
where $B = -(H_L + H_D)^{-1}H_U$.
\end{theorem}

\begin{proof}
We rewrite the recurrence relation of Gibbs updates derived in~\eqref{eq:5.3} as
\[
  X^{(n)} - \mu = B(X^{(n-1)} - \mu) + z^{(n)},
\]
where $z^{(n)} \overset{\mathrm{i.i.d.}}{\sim} \Normal(0, \Sigma_z)$ for some $\Sigma_z$ that
we will determine in what follows. Taking expectation on both sides gives
\[
  \Expect[X^{(n)} - \mu] = B\,\Expect[X^{(n-1)} - \mu].
\]
On the other hand, we know that $\Normal(\mu, \Sigma)$ is the invariant distribution, so
$\mathbb{V}[x] = \mathbb{V}[Bx] + \mathbb{V}[z]$, where $x$ has the target distribution and
$z$ is the noise in the recurrence equation. This implies that
$\Sigma_z = \Sigma - B\Sigma B^T$. Taking the variance on both sides of the original recurrence
relation yields
\begin{align*}
  \mathbb{V}[X^{(n)}]
    &= \mathbb{V}[BX^{(n-1)}] + \Sigma_z\\
    &= B\,\mathbb{V}[X^{(n-1)}]\,B^T + \Sigma - B\Sigma B^T\\
    &= \Sigma + B\Bigl(\mathbb{V}[X^{(n-1)}] - \Sigma\Bigr)B^T.
\end{align*}
Let $\mu_n$ and $\Sigma_n$ be the mean and variance of the conditional distribution
$X^{(n)} | x^{(0)}$. Then,
\begin{align*}
  \mu_n     &= \mu + B^n(x^{(0)} - \mu),\\
  \Sigma_n  &= \Sigma + B^n(0 - \Sigma)(B^T)^n = \Sigma - B^n\Sigma(B^T)^n,
\end{align*}
because $\Sigma^{(0)} = 0$. Using Equation~\eqref{eq:5.4}, we have
\[
  d_{\chi^2}\bigl(p_{\mathrm{DUGS}}^n(x^{(0)}, \cdot)\,\|\, f\bigr)
  = \frac{|W|}{\sqrt{|2W - I|}}\,e^{u^T(\Sigma^{(n)})^{-1}(2W-I)^{-1}u} - 1,
\]
with
\begin{align*}
  W &= \Sigma_n\Sigma^{-1} = I - B^n\Sigma(B^T)^n,\\
  u &= B^n(x^{(0)} - \mu).
\end{align*}
Observe that the only $x^{(0)}$-dependence lies in the exponent. Integrating over
$x^{(0)} \sim f$ gives
\begin{align*}
  \Expect_f\!\Bigl[d_{\chi^2}\bigl(p_{\mathrm{DUGS}}^n(x^{(0)}, \cdot)\,\|\, f\bigr)\Bigr] + 1
  &= \frac{|W|}{\sqrt{|2W - I|}}\int e^{u^T(\Sigma^{(n)})^{-1}(2W-I)^{-1}u} f(x)\,dx\\
  &= \frac{|W|}{\sqrt{|2W - I|}}\cdot\frac{1}{(2\pi)^{\frac{d}{2}}|\Sigma|^{\frac{1}{2}}}
    \int e^{\frac{1}{2}(x-\mu)^T(2(B^T)^t\Sigma^{-1}W^{-1}(2W-I)^{-1}B^t - \Sigma^{-1})(x-\mu)}\\
  &= \frac{|W|}{\sqrt{|2W - I|\,|U|}},
\end{align*}
where $U = I - 2(B^T)^n\Sigma^{-1}W^{-1}(2W-I)^{-1}B^n\Sigma$. From the definition of $W$,
we obtain the desired rate $\rho(B)$.
\end{proof}

\begin{figure}[htbp]
  \centering
  \includegraphics[width=\textwidth]{figures/fig_05_3}
  \caption{\textit{Gibbs samplers for a zero-mean bivariate Gaussian distribution with
    covariance matrix $[1, \alpha;\, \alpha, 1]$. Top row: $\alpha = 0.01$.
    Bottom row: $\alpha = 0.99$.}}
  \label{fig:5.3}
\end{figure}
